\documentclass[11pt,a4paper]{article}

% ─────────────────────────────────────────────────────────────────────────────
% A note for Matthew Colbrook on the functional-analytic setting of operator
% folding, written ahead of a video call.  Macro names deliberately match
% ~/Maths/Notes/spectral_theory/shared.tex so that material can migrate into
% the survey later.
% ─────────────────────────────────────────────────────────────────────────────

\usepackage[margin=3cm]{geometry}
\usepackage{amsmath}
\usepackage{amssymb}
\usepackage{amsfonts}
\usepackage{amsthm}
\usepackage{mathtools}
\usepackage{bm}
\usepackage{tikz}
\usepackage{tikz-cd}
\usepackage[colorlinks=true,linkcolor=blue!55!black,citecolor=blue!55!black,
            urlcolor=blue!55!black]{hyperref}
\usepackage{cleveref}
% thmtools is loaded after cleveref on purpose.  Environments declared with
% \declaretheorem register their own cleveref type even when they share a
% counter, which plain \newtheorem[theorem] does not.
\usepackage{thmtools}

\theoremstyle{plain}
\newtheorem{theorem}{Theorem}[section]
\declaretheorem[sibling=theorem,style=plain,name=Proposition]{proposition}
\declaretheorem[sibling=theorem,style=plain,name=Lemma]{lemma}
\declaretheorem[sibling=theorem,style=plain,name=Corollary]{corollary}
\declaretheorem[sibling=theorem,style=definition,name=Definition]{definition}
\declaretheorem[sibling=theorem,style=definition,name=Example]{example}
\declaretheorem[style=definition,name=Question]{question}
\declaretheorem[sibling=theorem,style=remark,name=Remark]{remark}

\newcommand{\creflastconjunction}{, and~}

% All of the theorem-like environments above share the `theorem' counter, so
% cleveref would otherwise call every one of them a theorem.
\crefname{definition}{Definition}{Definitions}
\Crefname{definition}{Definition}{Definitions}
\crefname{lemma}{Lemma}{Lemmas}
\Crefname{lemma}{Lemma}{Lemmas}
\crefname{proposition}{Proposition}{Propositions}
\Crefname{proposition}{Proposition}{Propositions}
\crefname{corollary}{Corollary}{Corollaries}
\Crefname{corollary}{Corollary}{Corollaries}
\crefname{remark}{Remark}{Remarks}
\Crefname{remark}{Remark}{Remarks}
\crefname{example}{Example}{Examples}
\Crefname{example}{Example}{Examples}
\crefname{theorem}{Theorem}{Theorems}
\Crefname{theorem}{Theorem}{Theorems}
\crefname{question}{Question}{Questions}
\Crefname{question}{Question}{Questions}

% ── Macros (names shared with the spectral_theory survey) ────────────────────
\newcommand{\R}{\mathbb{R}}
\newcommand{\C}{\mathbb{C}}
\newcommand{\N}{\mathbb{N}}
\newcommand{\norm}[1]{\left\|#1\right\|}
\newcommand{\abs}[1]{\left|#1\right|}
\newcommand{\dual}[2]{\left\langle #1, #2 \right\rangle}
\newcommand{\Hb}{\mathcal{H}}
\newcommand{\Vb}{V}
\newcommand{\Lop}{\mathcal{L}}
\newcommand{\Aop}{\mathcal{A}}
\newcommand{\Bop}{\mathcal{B}}
\newcommand{\Dop}{\mathcal{D}}
\newcommand{\Mop}{\mathcal{M}}
\newcommand{\spec}{\sigma}
\newcommand{\ii}{\mathrm{i}}
\newcommand{\dd}{\,\mathrm{d}}
\DeclareMathOperator{\dom}{dom}
\DeclareMathOperator{\Sp}{Sp}
\DeclareMathOperator{\dist}{dist}
\DeclareMathOperator{\Real}{Re}
\DeclareMathOperator{\rot}{rot}
\DeclareMathOperator{\Rot}{\mathbf{rot}}
\DeclareMathOperator{\curl}{curl}
\DeclareMathOperator{\divg}{div}
\newcommand{\siginf}{\spec_{\inf}}
\newcommand{\adj}{\star}                 % antidual (Banach) adjoint
\newcommand{\Riesz}{\mathcal{R}}         % the H-Riesz map
\newcommand{\Mass}{\mathcal{M}}          % the mass operator V -> V'
\newcommand{\Ltwo}{L^2}
\newcommand{\Hdiv}{H(\divg)}
\newcommand{\Hrot}{H(\rot)}
\newcommand{\HRot}{H(\Rot)}

\title{Operator folding and the primal-dual bookkeeping\\ of variational
       problems\\[2ex]
       \large A note for Matthew Colbrook}
\author{Umberto Zerbinati\thanks{Mathematical Institute, University of Oxford
  (\href{mailto:zerbinati@maths.ox.ac.uk}{zerbinati@maths.ox.ac.uk}).  Written
  after discussions with Margaret McCarthy and Patrick E. Farrell.}}
\date{\today}

\begin{document}
\maketitle

\begin{abstract}
Chapter~3 of \cite{Colbrook2026} builds \texttt{PseudoSpec} and \texttt{CompSpec}
on top of a closed and densely defined operator $A \colon \dom(A) \subset \Hb
\to \Hb$, an endomorphism of a single Hilbert space.  A finite element method
does not hand you such an object.  It hands you a bounded operator $\Aop \colon
\Vb \to \Vb'$ from a space into its own dual, which is an endomorphism of
nothing, and the passage from one picture to the other is where we kept getting
stuck.  This note works
out that passage.  The short answer is that there is no primal-dual
inconsistency in operator folding, there is a Riesz map which is invisible in
$\ell^2(\N)$ and unavoidable in $\Vb'$, and it is already present in
\cite[Algorithm~3.4]{Colbrook2026} under the name $G_n$.  Two further worries,
about Kato's domain for $(A-zI)^*(A-zI)$ and about the absence of a compact
resolvent, both dissolve once one notices that the algorithm touches only the
\emph{form} domain and only the \emph{bottom} of the spectrum.  We close with
four questions.
\end{abstract}

\tableofcontents

% ═════════════════════════════════════════════════════════════════════════════
\section{What this note is about}
\label{sec:purpose}

Maggie's thesis applies two layers of adaptivity to the algorithms of
\cite[Chapter~3]{Colbrook2026}.  An inner layer refines the discretisation used
to evaluate the injection moduli, and an outer layer refines the sampling of
the complex plane.  Both layers live on top of a finite element discretisation,
and it was while writing down what exactly is being discretised that we ran into
trouble.

Recall the two objects.  For a closed and densely defined $A$ on a separable
Hilbert space $\Hb$, the injection modulus is
\begin{equation}
  \label{eq:siginf}
  \siginf(A) = \inf\{\norm{Au} : u \in \dom(A),\ \norm{u} = 1\},
\end{equation}
and the resolvent lower bound that drives everything is
\begin{equation}
  \label{eq:gamma}
  \gamma(z,A) = \min\bigl\{\siginf(A - zI),\ \siginf(A^* - \bar z I)\bigr\}
  = \norm{(A-zI)^{-1}}^{-1},
\end{equation}
the last equality being \cite[Lemma~1.2.5]{Colbrook2026}.  Operator folding
\cite[\S3.4.1, p.~114]{Colbrook2026} evaluates \cref{eq:gamma} through the two
non-negative closed quadratic forms
\begin{equation}
  \label{eq:colbrookforms}
  Q^{(1)}_{z,A}(u,v) = \dual{(A-zI)u}{(A-zI)v},\qquad
  Q^{(2)}_{z,A}(u,v) = \dual{(A-zI)^*u}{(A-zI)^*v},
\end{equation}
with $\dom(Q^{(1)}_{z,A}) = \dom(A)$ and $\dom(Q^{(2)}_{z,A}) = \dom(A^*)$, so
that
\begin{equation}
  \label{eq:326}
  \gamma(z,A) = \sqrt{\min\Bigl\{
     \inf_{u \in \dom(A)\setminus\{0\}} \tfrac{Q^{(1)}_{z,A}(u,u)}{\norm{u}^2},\
     \inf_{u \in \dom(A^*)\setminus\{0\}} \tfrac{Q^{(2)}_{z,A}(u,u)}{\norm{u}^2}
  \Bigr\}}.
\end{equation}
This is \cite[Equation~(3.26)]{Colbrook2026}.  The sentence immediately after it
says that the computation of $\gamma$ has been transformed into finding the
bottom of the spectrum of the non-negative self-adjoint operators
$(A-zI)^*(A-zI)$ and $(A-zI)(A-zI)^*$, and it is that sentence which started our
confusion.

Everything in \cref{eq:siginf,eq:gamma,eq:326} presupposes that $A$ maps $\Hb$
into itself.  A finite element method starts from a bilinear form and produces
$\Aop \colon \Vb \to \Vb'$, so the composition $\Aop^{*}\Aop$ appears not to
type-check, and the standard cure of silently identifying $\Hb$ with $\Hb'$ is
exactly the sort of thing that hides a factor of a mass matrix.  Our four
questions were the following.

\begin{enumerate}
  \item How does one legitimately pass from $\Aop \colon \Vb \to \Vb'$ to
        Colbrook's $A \colon \dom(A) \subset \Hb \to \Hb$?
  \item Is there a primal-dual inconsistency in the folded operator, and if
        the cure is a Riesz map, which Riesz map?
  \item Kato \cite[Theorem~V.3.24, p.~275]{Kato1995} constructs $(A-zI)^*(A-zI)$
        on a domain that is strictly smaller than $\dom(A)$ and is not given
        explicitly.  Do we need to discretise that domain?
  \item The folded operator is unbounded and has no compact inverse, so
        Babu\v{s}ka--Osborn theory \cite{BabuskaOsborn1991,Boffi2010} does not
        apply to it.  What replaces it?
\end{enumerate}

The answers, in order, are as follows, and they occupy
\cref{sec:bridge,sec:folding,sec:gram,sec:kato,sec:rayleigh}.

\begin{itemize}
  \item[(1)] $A = \Riesz \Aop$ restricted to $\dom(A) = \{u \in \Vb : \Aop u
        \in \Hb'\}$, where $\Riesz \colon \Hb' \to \Hb$ is the Riesz map of the
        \emph{pivot} space.  Two things happen at once.  The Riesz map
        transports the output back into $\Hb$, and the domain shrinks from $\Vb$
        to $\Aop^{-1}(\Hb')$.  The second is the one that gets forgotten, and it
        is where the boundary conditions and the elliptic regularity live.
  \item[(2)] There is no inconsistency.  The Banach adjoint of $\Aop \colon \Hb
        \to \Hb'$ is a map $\Hb'' = \Hb \to \Hb'$ and not a map $\Hb' \to \Hb'$,
        so the two arrows simply do not compose and a Riesz map has to be
        inserted between them.  The correct folded object is $(\Aop -
        z\Mass)^{\adj}\Riesz(\Aop - z\Mass) \colon \Vb \to \Vb'$, and the Riesz
        map is the one of $\Hb$, not the one of $\Vb$.  This is the continuous
        counterpart of the modified normal equation of
        \cite[Equation~(3.7)]{LazzarinoNakatsukasaZerbinati2026}.
  \item[(3)] No.  Rayleigh quotients see only the \emph{form} domain, which by
        \cite[Example~VI.2.13, p.~326]{Kato1995} is exactly $\dom(A)$.  Kato's
        smaller domain is a core of that form and never has to be exhibited.
  \item[(4)] Nothing needs to replace it.  The folded form is non-negative and
        only its infimum is wanted, so a Rayleigh--Ritz restriction is
        automatically an upper bound.  That is precisely the one-sidedness
        $\gamma(z,A) \le \Phi_n(z,A)$ demanded by
        \cite[Assumption~3.2.1]{Colbrook2026}.  The approximation hypothesis is
        that the trial spaces exhaust a \emph{core}, which is density in the
        graph norm and not compactness.
\end{itemize}

\Cref{sec:twodomains,sec:tworoutes} record two consequences that are invisible
in $\ell^2(\N)$ and very visible in a finite element code.  \Cref{sec:examples}
carries the argument through Maxwell, Dirac and advection-diffusion.
\Cref{sec:affine} records a cheap way to sweep the complex plane.
\Cref{sec:pivot} asks how far the pivot space can be moved, which is where this
touches the preconditioning work in
\cite{LazzarinoNakatsukasaZerbinati2026}.  \Cref{sec:questions} collects the
questions we would like to discuss.

% ═════════════════════════════════════════════════════════════════════════════
\section{Notation, fixed once}
\label{sec:notation}

Sesquilinear forms and inner products are linear in the first argument and
conjugate-linear in the second.  We use \emph{antiduals} throughout.  For a
complex Hilbert space $X$, $X'$ denotes the space of bounded conjugate-linear
functionals on $X$, and $\dual{g}{v} := g(v)$ for $g \in X'$ and $v \in X$, so
that $\dual{\cdot}{\cdot}$ is linear in its first slot and conjugate-linear in
its second.  With this convention the Riesz map is linear rather than
conjugate-linear, and no bars appear where they are not wanted.  This is a
bookkeeping choice and nothing depends on it, but it removes a genuine share of
the confusion.

Let $\Hb$ be the \emph{pivot} Hilbert space, typically $L^2$, and let $\Vb
\hookrightarrow \Hb$ be a Hilbert space that is continuously and densely
embedded.  We write
\begin{equation}
  \label{eq:riesz}
  \Riesz \colon \Hb' \to \Hb, \qquad \dual{g}{w} = (\Riesz g, w)_{\Hb}
  \quad \text{for all } w \in \Hb,
\end{equation}
for the Riesz map of $\Hb$, which is a linear isometric isomorphism, and
\begin{equation}
  \label{eq:mass}
  \Mass \colon \Vb \to \Vb', \qquad \dual{\Mass u}{v} = (u,v)_{\Hb}
  \quad \text{for all } v \in \Vb,
\end{equation}
for the \emph{mass operator}, which is the composition of the inclusion $\Vb
\hookrightarrow \Hb$, the inverse Riesz map $\Hb \to \Hb'$, and the restriction
of functionals $\Hb' \hookrightarrow \Vb'$.  The last of these is injective
because $\Vb$ is dense in $\Hb$.  The Gelfand triple is
\begin{equation}
  \label{eq:gelfand}
  \begin{tikzcd}[column sep=large]
    \Vb \arrow[r, hook] & \Hb \arrow[r, "\Riesz^{-1}", "\sim"'] & \Hb'
    \arrow[r, hook] & \Vb',
  \end{tikzcd}
\end{equation}
and we never write $\Hb = \Hb'$.  Two identities are used constantly,
\begin{equation}
  \label{eq:RM}
  \Riesz \Mass = \iota \colon \Vb \hookrightarrow \Hb, \qquad
  \norm{\Riesz g}_{\Hb} = \norm{g}_{\Hb'} \quad \text{for } g \in \Hb'.
\end{equation}

Given a bounded operator $\mathcal{C} \colon \Vb \to \Hb'$ we write
$\mathcal{C}^{\adj} \colon \Hb \to \Vb'$ for its antidual adjoint,
\begin{equation}
  \label{eq:antidualadjoint}
  \dual{\mathcal{C}^{\adj} w}{v} := \overline{\dual{\mathcal{C}v}{w}},
  \qquad w \in \Hb,\ v \in \Vb.
\end{equation}
The right-hand side is conjugate-linear in $v$ and linear in $w$, so
$\mathcal{C}^{\adj}$ is a well-defined bounded linear map with
$\norm{\mathcal{C}^{\adj}} = \norm{\mathcal{C}}$.  Note the shape of the
arrows.  The adjoint of a map \emph{into} a dual space is a map \emph{out of}
the predual, not a map between duals.  This is the whole of \cref{sec:folding}
in one sentence.

Finally, $\Aop \in \mathcal{L}(\Vb,\Vb')$ denotes a bounded weak operator, with
associated sesquilinear form $a(u,v) := \dual{\Aop u}{v}$, and $\Aop^{\adj} \in
\mathcal{L}(\Vb,\Vb')$ denotes the operator of the adjoint form,
$\dual{\Aop^{\adj} v}{u} := \overline{a(u,v)}$.

% ═════════════════════════════════════════════════════════════════════════════
\section{From a variational operator to Colbrook's endomorphism}
\label{sec:bridge}

\begin{definition}[the strong realisation]
  \label{def:strong}
  Given $\Aop \in \mathcal{L}(\Vb,\Vb')$, set
  \begin{equation}
    \label{eq:domA}
    \dom(A) := \{u \in \Vb : \Aop u \in \Hb'\},
    \qquad A := \Riesz \Aop \big|_{\dom(A)} .
  \end{equation}
\end{definition}

Unwinding \cref{eq:riesz,eq:mass}, \cref{def:strong} says exactly
\begin{equation}
  \label{eq:strongchar}
  u \in \dom(A) \text{ and } Au = f
  \quad \Longleftrightarrow \quad
  u \in \Vb,\ f \in \Hb, \text{ and } a(u,v) = (f,v)_{\Hb}
  \ \ \forall v \in \Vb ,
\end{equation}
which is the usual recipe for reading a strong operator off a weak form.  Two
separate things are happening and it is worth separating them.  The Riesz map
transports the output of $\Aop$ from $\Hb'$ back into $\Hb$.  Independently, the
domain \emph{shrinks} from $\Vb$ to $\Aop^{-1}(\Hb')$, which is in general a
proper subspace.  For $\Aop = -\Delta$ with $\Vb = H^1_0(\Omega)$ and $\Hb =
L^2(\Omega)$ one gets $\dom(A) = H^2(\Omega) \cap H^1_0(\Omega)$ on a convex or
$C^{1,1}$ domain, strictly inside $\Vb$.  The shrinking is where the elliptic
regularity and the boundary conditions live, and it is the step that gets
forgotten.  In the language of \cref{sec:purpose}, the endomorphism $A$ that
\cite{Colbrook2026} needs is not obtained from $\Aop$ by an identification, it
is obtained by a composition and a restriction.

Whether $A$ is closed and densely defined has to be argued, and the argument
differs from example to example.

\begin{remark}[provenance of \cref{def:strong}]
  \label{rem:provenance}
  If the form $a$ is densely defined, closed and \emph{sectorial}, then Kato's
  first representation theorem \cite[Theorem~VI.2.1, p.~322]{Kato1995} produces
  an m-sectorial, hence closed and densely defined, operator, and part~(iii) of
  that theorem characterises its domain exactly as the right-hand side of
  \cref{eq:domA}.  So \cref{def:strong} \emph{is} Kato's construction, written
  with the Riesz map made explicit.  Sectoriality is a real hypothesis and it
  fails for indefinite first-order systems.  For Maxwell in \cref{sec:maxwell}
  the form is symmetric but bounded neither above nor below, and there one does
  not go through \cref{def:strong} at all.  Instead the self-adjoint operator is
  written down directly with $\dom(A) = \Vb$, as in
  \cite[p.~130]{Colbrook2026}.
\end{remark}

The adjoint behaves as one would hope, with the same caveat.

\begin{lemma}[the three adjoints agree, up to Riesz maps]
  \label{lem:adjoint}
  Let $\Aop \in \mathcal{L}(\Vb,\Vb')$ and let $A$ be as in
  \cref{def:strong}.  Put $\dom(A^{\adj}) := \{v \in \Vb : \Aop^{\adj} v \in
  \Hb'\}$ and $A^{\adj}_{\Riesz} := \Riesz\Aop^{\adj}|_{\dom(A^{\adj})}$.  Then
  \begin{equation}
    \label{eq:adjinclusion}
    A^{\adj}_{\Riesz} \subseteq A^*,
  \end{equation}
  where $A^*$ is the Hilbert adjoint of $A$ in $\Hb$.  Equality holds whenever
  $\dom(A^*) \subseteq \Vb$, in particular for closed sectorial forms.
\end{lemma}

\begin{proof}
  Let $v \in \dom(A^{\adj})$ and put $g := \Riesz\Aop^{\adj}v \in \Hb$.  For
  $u \in \dom(A) \subseteq \Vb$,
  \[
    (Au,v)_{\Hb} = \dual{\Aop u}{v} = a(u,v)
      = \overline{\dual{\Aop^{\adj}v}{u}}
      = \overline{(g,u)_{\Hb}} = (u,g)_{\Hb},
  \]
  using \cref{eq:riesz} and the definition of $\Aop^{\adj}$.  Hence $v \in
  \dom(A^*)$ with $A^*v = g$, which is \cref{eq:adjinclusion}.  The converse
  inclusion needs $\dom(A^*) \subseteq \Vb$, after which the same computation
  run backwards shows $\Aop^{\adj}v \in \Hb'$.
\end{proof}

So the Banach adjoint of the weak operator, composed with the Riesz map of the
pivot space, is the Hilbert adjoint of the strong operator.  The adjoint with
respect to the duality pairing and the adjoint with respect to the inner product
are not in competition.  They differ by Riesz maps which are explicit, and
\cref{eq:adjinclusion} is the precise statement of that.  This answers a
question we had left open in an appendix of our own working notes.

The last consequence of \cref{def:strong} is the one that fixes the shift.

\begin{corollary}[the correct weak shift is by the mass operator]
  \label{cor:shift}
  For $u \in \dom(A)$ and $z \in \C$,
  \begin{equation}
    \label{eq:shift}
    (A - zI)u = \Riesz(\Aop - z\Mass)u .
  \end{equation}
\end{corollary}

\begin{proof}
  $zu = z\Riesz\Mass u$ by \cref{eq:RM}, and $\Riesz$ is linear.
\end{proof}

% ═════════════════════════════════════════════════════════════════════════════
\section{Folding without a primal-dual error}
\label{sec:folding}

Write $\mathcal{C}_z := \Aop - z\Mass$.  By \cref{eq:domA} we have
$\mathcal{C}_z(\dom(A)) \subseteq \Hb'$, and if $\dom(A) = \Vb$ then
$\mathcal{C}_z \colon \Vb \to \Hb'$ is bounded.

The diagram which does not close is
\begin{equation}
  \label{eq:baddiagram}
  \begin{tikzcd}[column sep=huge]
    \Hb \arrow[r, "\Aop"] & \Hb'
    & \Hb' \arrow[r, "\Aop^{\adj}", dashed] & \Hb' ,
  \end{tikzcd}
\end{equation}
and it fails for a reason that can be stated in one line.  The antidual adjoint
of $\Aop \colon \Hb \to \Hb'$ is a map $\Hb'' = \Hb \to \Hb'$, not a map
$\Hb' \to \Hb'$, by \cref{eq:antidualadjoint}.  The two arrows are simply not
composable, and the missing arrow between them is the Riesz map.  Inserting it
gives the diagram which does close,
\begin{equation}
  \label{eq:gooddiagram}
  \begin{tikzcd}[column sep=huge]
    \Vb \arrow[r, "\mathcal{C}_z"] & \Hb' \arrow[r, "\Riesz", "\sim"']
    & \Hb \arrow[r, "\mathcal{C}_z^{\adj}"] & \Vb' ,
  \end{tikzcd}
\end{equation}
so that the \emph{folded operator} is
\begin{equation}
  \label{eq:folded}
  \boxed{\ \Bop_z := (\Aop - z\Mass)^{\adj}\,\Riesz\,(\Aop - z\Mass)
    \ \colon\ \Vb \longrightarrow \Vb' .\ }
\end{equation}
Nothing has been identified and nothing has been lost.  In particular the
symbol $\Mass^{-1}$ that one is tempted to write in the middle of
\cref{eq:folded} means $\Riesz$, the Riesz map of the \emph{pivot} space $\Hb$,
and not the inverse of $\Mass \colon \Vb \to \Vb'$, which does not exist as a
map $\Vb' \to \Vb$.  The composition is legitimate precisely because
$\mathcal{C}_z$ lands in $\Hb' \subsetneq \Vb'$, which is the same statement as
$\dom(A) = \Vb$.

The bridge between \cref{eq:folded} and \cref{eq:colbrookforms} is the following
identity, which is the single most useful line of this note.

\begin{proposition}[the residual identity]
  \label{prop:residual}
  For $u \in \dom(A)$ and $z \in \C$,
  \begin{equation}
    \label{eq:residual}
    \norm{(A - zI)u}_{\Hb} = \norm{(\Aop - z\Mass)u}_{\Hb'} ,
  \end{equation}
  and consequently, for $u,v \in \dom(A)$,
  \begin{equation}
    \label{eq:formfolded}
    Q^{(1)}_{z,A}(u,v)
      = \bigl((A-zI)u,\ (A-zI)v\bigr)_{\Hb}
      = \dual{\Bop_z u}{v}_{\Vb' \times \Vb} .
  \end{equation}
\end{proposition}

\begin{proof}
  \Cref{eq:residual} is \cref{eq:shift} together with the isometry in
  \cref{eq:RM}.  For \cref{eq:formfolded}, use \cref{eq:antidualadjoint} and
  then \cref{eq:riesz},
  \[
    \dual{\mathcal{C}_z^{\adj}\Riesz\mathcal{C}_z u}{v}
      = \overline{\dual{\mathcal{C}_z v}{\Riesz \mathcal{C}_z u}}
      = \overline{(\Riesz\mathcal{C}_z v,\ \Riesz \mathcal{C}_z u)_{\Hb}}
      = (\Riesz\mathcal{C}_z u,\ \Riesz \mathcal{C}_z v)_{\Hb},
  \]
  and apply \cref{eq:shift} to both slots.
\end{proof}

In words, the $\Hb$-norm of the strong residual is the $\Hb'$-norm of the weak
residual.  The Riesz map is an isometry, so passing between the two costs
nothing, and the folded operator \cref{eq:folded} is the variational object
whose Rayleigh quotient is the first quotient of \cref{eq:326}.  The generalised
eigenvalue problem to be solved is
\begin{equation}
  \label{eq:pencil}
  \Bop_z u = \lambda\, \Mass u \quad \text{in } \Vb' ,
  \qquad \siginf(A - zI)^2 = \min \Sp(\Bop_z, \Mass),
\end{equation}
with the mass operator on the right-hand side and not the identity.  The second
quotient of \cref{eq:326} is the same construction run with $\Aop^{\adj}$ and
$\bar z$, by \cref{lem:adjoint}.

\begin{remark}[relation to the modified normal equation]
  \label{rem:lnz}
  \Cref{eq:folded} is the continuous counterpart of the modified normal equation
  $A^{T}TAx = A^{T}Tb$ of
  \cite[Equation~(3.7)]{LazzarinoNakatsukasaZerbinati2026}, in which $T$
  discretises a Riesz map $\tau \colon \Vb_h' \to \Vb_h$ inserted for exactly
  the reason given above, that the plain normal equation performs operations on
  inconsistent function spaces.  The point of contact, and also the point of
  difference, is \emph{which} Riesz map.  Folding forces $\tau = \Riesz$, the
  Riesz map of the pivot space, because that is the space in which
  $\siginf$ is measured.  We return to this in \cref{sec:pivot}.
\end{remark}

% ═════════════════════════════════════════════════════════════════════════════
\section{Where the Riesz map went in Algorithm 3.4}
\label{sec:gram}

It is worth saying plainly that none of this is a correction.  It is already in
\cite[Algorithm~3.4, p.~115]{Colbrook2026}.  That algorithm takes linearly
independent vectors $\{e_j\}$, which are explicitly \emph{not} assumed
orthonormal, and forms the three matrices
\begin{equation}
  \label{eq:alg34}
  \bigl[\mathsf{L}^{(1)}(z)\bigr]_{ij} = \dual{(A-zI)e_j}{(A-zI)e_i},\quad
  \bigl[\mathsf{L}^{(2)}(z)\bigr]_{ij} = \dual{(A-zI)^*e_j}{(A-zI)^*e_i},\quad
  \bigl[\mathsf{G}\bigr]_{ij} = \dual{e_j}{e_i},
\end{equation}
and then minimises the \emph{generalised} Rayleigh quotients $x^*
\mathsf{L}^{(j)}(z)x / x^*\mathsf{G}x$.  Writing $\mathsf{A}_{ij} = a(e_j,e_i)$
and $\mathsf{M}_{ij} = (e_j,e_i)_{\Hb}$ for the stiffness and mass matrices of a
finite element basis, \cref{eq:formfolded} says
\begin{equation}
  \label{eq:matrixfolded}
  \mathsf{L}^{(1)}(z)
    = (\mathsf{A} - z\mathsf{M})^{*}\,\mathsf{M}^{-1}\,(\mathsf{A} - z\mathsf{M}),
  \qquad \mathsf{G} = \mathsf{M} ,
\end{equation}
so that Colbrook's Gram matrix \emph{is} the mass matrix and his generalised
eigenvalue problem \emph{is} the discretisation of the pencil
\cref{eq:pencil}.  His own remark that ``the matrix $G_n$ is crucial for
capturing the correct norm in the computation'' \cite[p.~115]{Colbrook2026} is
this statement.  The framework is primal-dual consistent as it stands, and the
consistency is invisible only when $\Hb = \ell^2(\N)$ and the basis is
orthonormal, in which case $\mathsf{M} = I$ and all the Riesz maps disappear.

There is a second, practical observation hiding in \cref{eq:matrixfolded}.  The
matrix $\mathsf{M}^{-1}$ never has to be formed.  Folding evaluates
$\int_\Omega \abs{(A - zI)u_h}^2$ by quadrature, directly from the strong
residual, and that is exact whenever $u_h$ lies in $\dom(A)$ so that $(A -
zI)e_j$ is an explicit piecewise polynomial.  Equivalently, if $W_h \supseteq
(A - zI)V_h$ is a computable finite dimensional space, then
\begin{equation}
  \label{eq:sparsefactor}
  \mathsf{L}^{(1)}(z) = \mathsf{D}(z)^{*}\,\mathsf{M}_{W}\,\mathsf{D}(z),
\end{equation}
where $\mathsf{D}(z)$ represents $(A - zI) \colon V_h \to W_h$ and is sparse,
and $\mathsf{M}_W$ is the mass matrix of $W_h$, which is block diagonal if $W_h$
is discontinuous.  This matters because the dense factor $\mathsf{M}^{-1}$ in
\cref{eq:matrixfolded} is exactly the obstacle flagged in
\cite[Remark~3.6]{LazzarinoNakatsukasaZerbinati2026}, and
\cref{eq:sparsefactor} is how folding avoids it.

% ═════════════════════════════════════════════════════════════════════════════
\section{No special domain is needed}
\label{sec:kato}

Our third worry was that $(A-zI)^*(A-zI)$ is not well posed as a closed and
densely defined operator on $\dom(A)$, and that Kato constructs it on a domain
which is smaller and which is not given explicitly.  That is correct, and it is
also irrelevant, because the algorithm never touches that domain.

\begin{theorem}[Kato]
  \label{thm:kato}
  Let $S$ be a closed and densely defined operator from $\Hb$ to a Hilbert space
  $\Hb'$.  Then the following hold.
  \begin{enumerate}
    \item[(i)] The form $\mathfrak{h}[u,v] := (Su,Sv)$ with
      $\dom(\mathfrak{h}) = \dom(S)$ is densely defined, closed and
      non-negative, and the self-adjoint operator associated with it by the
      first representation theorem is exactly $S^*S$
      \cite[Example~VI.2.13, p.~326]{Kato1995}.
    \item[(ii)] $S^*S$ is self-adjoint and $\dom(S^*S)$ is a \emph{core} of $S$
      \cite[Theorem~V.3.24, p.~275]{Kato1995}.
    \item[(iii)] $\dom\bigl((S^*S)^{1/2}\bigr) = \dom(S)$ and
      $\norm{(S^*S)^{1/2}u} = \norm{Su}$ for $u \in \dom(S)$
      \cite[Theorem~VI.2.23 and Equation~(VI.2.22), pp.~331 and 334]{Kato1995}.
  \end{enumerate}
\end{theorem}

Apply this with $S = A - zI$, which is closed and densely defined because $A$
is.  Part~(i) says that the \emph{form} domain of the folded operator is
$\dom(A)$, no more and no less, and part~(iii) says that the form is the square
of the norm of $\abs{A - zI} = ((A-zI)^*(A-zI))^{1/2}$ on that same domain.  The
mysterious smaller domain of part~(ii) is $\dom(S^*S)$, and it is a core, so
\begin{equation}
  \label{eq:coreirrelevant}
  \inf_{u \in \dom(S^*S)\setminus\{0\}} \frac{(S^*Su,u)}{\norm{u}^2}
  = \inf_{u \in \dom(S)\setminus\{0\}} \frac{\norm{Su}^2}{\norm{u}^2}
  = \siginf(S)^2 .
\end{equation}
Both infima are the bottom of $\Sp(S^*S)$ and the left-hand one is never the one
that is computed.  This is why \cite[\S3.4.1]{Colbrook2026} is phrased in terms
of quadratic \emph{forms} rather than operators, and why the sentence after
\cite[Equation~(3.26)]{Colbrook2026} about the spectrum of $(A-zI)^*(A-zI)$
carries no hidden domain obligation.  Kato's second representation theorem is
what licenses the translation between the two phrasings.

% ═════════════════════════════════════════════════════════════════════════════
\section{One-sided Rayleigh--Ritz replaces Babu\v{s}ka--Osborn}
\label{sec:rayleigh}

Our fourth worry was that the folded operator is unbounded with non-compact
inverse, so that the Babu\v{s}ka--Osborn theory of \cite{BabuskaOsborn1991} and
the framework surveyed in \cite{Boffi2010} do not apply.  They do not, and they
are not needed.  Two features of the problem conspire to make a much weaker
theory sufficient.

First, $\Bop_z \ge 0$ and only the \emph{bottom} of its spectrum is wanted.  For
any subspace $\Dop_h \subseteq \dom(A)$,
\begin{equation}
  \label{eq:onesided}
  \Phi_h(z)^2 := \min_{u \in \Dop_h \setminus\{0\}}
     \frac{\norm{(A-zI)u\vphantom{|}}^2_{\Hb}}{\norm{u}^2_{\Hb}}
  \ \ge\ \siginf(A - zI)^2 ,
\end{equation}
simply because a minimum over a subset is at least the infimum over the whole
set.  This is the one-sidedness $\gamma(z,A) \le \Phi_n(z,A)$ that
\cite[Assumption~3.2.1]{Colbrook2026} demands, and which
\cite[Theorem~3.2.3]{Colbrook2026} turns into the inclusion $\Gamma_n^\epsilon(A)
\subseteq \Sp_\epsilon(A)$.  It costs nothing beyond conformity, $\Dop_h
\subseteq \dom(A)$.

Second, convergence needs only that the trial spaces exhaust a core.

\begin{proposition}[the approximation hypothesis is a core condition]
  \label{prop:core}
  Let $\Dop_1 \subseteq \Dop_2 \subseteq \cdots \subseteq \dom(A)$ with
  $\bigcup_h \Dop_h$ a core of $A$, that is, dense in $\dom(A)$ for the graph
  norm $\norm{u}^2_{\Hb} + \norm{Au}^2_{\Hb}$.  Then $\Phi_h(z) \downarrow
  \siginf(A - zI)$ as $h \to 0$, monotonically and locally uniformly in $z$.
\end{proposition}

The proof is the one given for \cite[Lemma~1.2.6]{Colbrook2026}, and the
hypothesis is exactly the one stated as an input to
\cite[Algorithm~3.4]{Colbrook2026}.  Three things are worth emphasising.

\begin{itemize}
  \item No compactness of any resolvent is used.
  \item The bottom of $\Sp(\Bop_z)$ is allowed to be \emph{essential} spectrum.
        This is precisely the case that Babu\v{s}ka--Osborn theory cannot treat
        and that a one-sided Rayleigh--Ritz bound treats without noticing.
  \item The density required is in the \emph{graph} norm, not in the energy
        norm of $\Vb$.  The right discrete space is a subspace of the graph
        space $\dom(A)$.
\end{itemize}

That last point is the practical content of this section.  If $\dom(A) = \Vb$,
as happens for first-order systems whose natural energy space is their graph
space, then any $\Vb$-conforming element is admissible and folding is a standard
conforming Galerkin method.  This is why \cite[\S3.4.4]{Colbrook2026} can use
plain Lagrange elements for Maxwell and observe no spectral pollution.  If
$\dom(A) \subsetneq \Vb$, as happens for a second-order operator with $\Vb =
H^1_0$, then $\Vb$-conforming elements are \emph{not} $\dom(A)$-conforming,
$\bigcup_h \Dop_h$ is not a core, and neither \cref{eq:onesided} nor
\cref{prop:core} is available.  \Cref{sec:advdiff} discusses what to do.

\begin{remark}[where compactness does belong]
  \label{rem:whereboffi}
  Compactness is not needed for the \emph{guarantee}, but it is exactly what is
  needed for a \emph{rate} and for an a posteriori \emph{gap}.  The
  superconvergence estimate \cite[Proposition~3.3.12, p.~107]{Colbrook2026} and
  the Kato--Temple inequality \cite[Theorem~3.3.11]{Colbrook2026} behind it
  require $\lambda$ to be an isolated point of the spectrum and a computable
  lower bound on the distance to the rest of it.  If the solution operator
  associated with $\Bop_z$ maps into a space compactly embedded in $\Vb$, then
  the bottom of $\Sp(\Bop_z)$ is an isolated eigenvalue and Boffi-style theory
  \cite{Boffi2010,BoffiBrezziFortin2013} delivers both.  The two frameworks are
  complementary rather than competing.  The core condition gives the certificate
  and compactness, when it is available, gives the speed.
\end{remark}

% ═════════════════════════════════════════════════════════════════════════════
\section{Two domains, not one}
\label{sec:twodomains}

Definition \cref{eq:326} involves both $\dom(A)$ and $\dom(A^*)$, and by
\cref{lem:adjoint} these are
\begin{equation}
  \label{eq:twodomains}
  \dom(A) = \{u \in \Vb : \Aop u \in \Hb'\},
  \qquad
  \dom(A^*) = \{v \in \Vb : \Aop^{\adj} v \in \Hb'\},
\end{equation}
which are different subspaces of $\Vb$ in general.  Consequently the two
Rayleigh quotients of \cref{eq:326} need \emph{two} trial spaces, one a core of
$A$ and the other a core of $A^*$, and \cref{eq:onesided} has to hold for each
separately.  In $\ell^2(\N)$ with a single orthonormal basis this is invisible,
and \cite[Algorithm~3.4]{Colbrook2026} accordingly asks for one family $\{e_j\}$
whose span is a core of both.  In a finite element code the two spaces usually
differ in their essential boundary conditions, and the distinction is very
visible.  For the two-dimensional Maxwell system of \cref{sec:maxwell} the two
spaces are $H_0(\rot) \times H(\Rot)$ and its swap, so a single choice works
only because that operator is self-adjoint.  For a genuinely non-normal problem
they must be built separately.

This is also where the boundary terms of
\cite[Equations~(3.5)--(3.6)]{LazzarinoNakatsukasaZerbinati2026} reappear.  The
adjoint form is obtained by integration by parts, and whether the resulting
boundary contributions vanish is a statement about the boundary conditions
encoded in $\Vb$, not about the differential expression.

% ═════════════════════════════════════════════════════════════════════════════
\section{Reading off a variational formulation, two routes}
\label{sec:tworoutes}

A recurring practical question was how to read a variational formulation off
$(\Aop - z\Mass)^{\adj}\Riesz(\Aop - z\Mass)$, given that a finite element code
wants a bilinear form and not a composition of maps.  By
\cref{prop:residual} the answer is that the bilinear form is
\begin{equation}
  \label{eq:qz}
  q_z(u,v) = \bigl((A - zI)u,\ (A - zI)v\bigr)_{\Hb},
  \qquad u,v \in \dom(A),
\end{equation}
and there are two ways to realise it.

\paragraph{Route (a), least squares.}  Keep $A$ in strong form and integrate the
product of the two residuals, $\int_\Omega (Au - zu)\,\overline{(Av - zv)}$.
This needs $V_h \subseteq \dom(A)$, it never forms $\mathsf{M}^{-1}$, and it
preserves the one-sidedness \cref{eq:onesided} exactly.  It is a first-order
system least squares formulation \cite{CaiLazarovManteuffelMcCormick1994} for
the shifted operator, and it is what \cite[Algorithm~3.4]{Colbrook2026} does.

\paragraph{Route (b), integrate by parts once.}  Move a derivative from the
first factor to the second.  This lowers the differential order of the form,
which is what makes conforming discretisation possible in some cases, at the
price of producing boundary terms.  Those boundary terms are precisely the ones
that encode $\dom(A^*)$, so route (b) is where the primal-dual bookkeeping of
\cref{sec:twodomains} has to be done honestly.  This is the manoeuvre referred
to in \cite[p.~129]{Colbrook2026} as reformulating ``into a mixed form using
lower-order derivatives''.

The two routes give the same number whenever the boundary terms of route (b)
vanish.  Route (a) is the robust one.  Route (b) is the one that produces the
bilinear form a finite element code expects, and \cref{sec:maxwell,sec:dirac}
carry it out.

% ═════════════════════════════════════════════════════════════════════════════
\section{Three examples}
\label{sec:examples}

\subsection{Maxwell, where everything works}
\label{sec:maxwell}

Following \cite[\S3.4.4, p.~130]{Colbrook2026}, take the two-dimensional
first-order Maxwell system on a bounded Lipschitz domain $U \subset \R^2$.  With
$\rot E = \partial_x E_2 - \partial_y E_1$ for vector fields and $\Rot H =
[\partial_y H, -\partial_x H]$ for scalars, the operator is
\begin{equation}
  \label{eq:maxwellop}
  A(E,H) = \bigl(\ii\,\Rot H,\ -\ii\,\rot E\bigr),
  \qquad
  \dom(A) = H_0(\rot;U) \times H(\Rot;U) =: \Vb ,
\end{equation}
on $\Hb = L^2(U)^2 \times L^2(U)$.  It is self-adjoint, and $\dom(A) = \Vb$ is
its graph space, so by \cref{sec:rayleigh} any $\Vb$-conforming element is
admissible.  This is the reason \cite{Colbrook2026} can use plain Lagrange
elements for both components, imposing the boundary condition only on the space
approximating $E$, and see no spectral pollution where the standard variational
formulation pollutes badly.

Expanding \cref{eq:qz} for \cref{eq:maxwellop} and using $\ii\overline{\ii} = 1$
gives, for $(E,H)$ and $(F,G)$ in $\Vb$,
\begin{multline}
  \label{eq:maxwellqz}
  q_z\bigl((E,H),(F,G)\bigr)
    = (\rot E, \rot F) + (\Rot H, \Rot G)
      + \abs{z}^2\bigl[(E,F) + (H,G)\bigr] \\
      {} + \ii\bar z\bigl[(\rot E, G) - (\Rot H, F)\bigr]
      + \ii z\bigl[(E, \Rot G) - (H, \rot F)\bigr],
\end{multline}
all inner products being those of $L^2(U)$.  This is route (a), and the two
shift terms are Hermitian conjugates of one another as they must be.  Now apply
route (b).  For $E \in H_0(\rot;U)$ and $G \in H(\Rot;U)$ the integration by
parts
\begin{equation}
  \label{eq:rotibp}
  (\rot E, G)_{L^2} = (E, \Rot G)_{L^2}
\end{equation}
holds with no boundary contribution, because the tangential trace of $E$
vanishes.  Using \cref{eq:rotibp} twice, once as written and once with the roles
of $(E,G)$ and $(F,H)$ exchanged, \cref{eq:maxwellqz} collapses to
\begin{multline}
  \label{eq:maxwellqzb}
  q_z\bigl((E,H),(F,G)\bigr)
    = (\rot E, \rot F) + (\Rot H, \Rot G)
      + \abs{z}^2\bigl[(E,F) + (H,G)\bigr] \\
      {} + \ii(z + \bar z)\bigl[(\rot E, G) - (\Rot H, F)\bigr] ,
\end{multline}
with $z + \bar z = 2\Real z$.  \Cref{eq:maxwellqzb} is the form we had already
written down by hand, and \cref{eq:maxwellqz} is Colbrook's minimisation
functional on \cite[p.~132]{Colbrook2026}.  They agree, and the agreement uses
\cref{eq:rotibp}, which is exactly the statement that $\rot$ on $H_0(\rot)$ and
$\Rot$ on $H(\Rot)$ are adjoint to one another.

This is worth pausing on, because it answers the question we set out with.  The
primal-dual identification in the folded Maxwell form is handled correctly, and
the reason it is handled correctly is the \emph{asymmetry} of the boundary
conditions between the two curl spaces.  If both components carried the
vanishing tangential trace, or neither did, \cref{eq:rotibp} would acquire a
boundary term and \cref{eq:maxwellqzb} would be wrong while
\cref{eq:maxwellqz} would remain right.  Route (a) is insensitive to this and
route (b) is not.

\subsection{Dirac, where the form is first order although the operator is not}
\label{sec:dirac}

Let $\alpha_1,\alpha_2,\alpha_3,\beta$ be the Hermitian $4 \times 4$ Dirac
matrices, so that $\alpha_j\alpha_k + \alpha_k\alpha_j = 2\delta_{jk}I$,
$\alpha_j\beta + \beta\alpha_j = 0$ and $\beta^2 = I$
\cite{Thaller1992}.  The free Dirac operator on $\Hb = L^2(\R^3;\C^4)$
is
\begin{equation}
  \label{eq:diracop}
  \Dop = -\ii c\, \alpha \cdot \nabla + mc^2 \beta,
  \qquad \dom(\Dop) = H^1(\R^3;\C^4),
\end{equation}
and it is self-adjoint.  Its square is the Klein--Gordon operator, $\Dop^2 =
(-c^2\Delta + m^2c^4)I_4$, which is second order.  Nevertheless the folded form
is first order.

\begin{lemma}
  \label{lem:diracform}
  For $u,v \in H^1(\R^3;\C^4)$,
  \begin{equation}
    \label{eq:diracbilinear}
    (\Dop u, \Dop v)_{L^2} = c^2(\nabla u, \nabla v)_{L^2}
      + m^2c^4 (u,v)_{L^2} .
  \end{equation}
\end{lemma}

\begin{proof}
  The two cross terms are $-\ii m c^3\bigl[(\alpha\cdot\nabla u, \beta v) -
  (\beta u, \alpha\cdot\nabla v)\bigr]$.  Integrating by parts in the first
  summand and moving $\beta$ across in the second gives, for each $j$,
  $-(u, \alpha_j\beta\,\partial_j v) - (u, \beta\alpha_j\,\partial_j v)
   = -(u, (\alpha_j\beta + \beta\alpha_j)\partial_j v) = 0$.
  For the leading term, Plancherel and $(\alpha\cdot\xi)^2 = \abs{\xi}^2 I$ give
  $(\alpha\cdot\nabla u, \alpha\cdot\nabla v) = (\nabla u, \nabla v)$.  The
  remaining term is $m^2c^4(\beta u, \beta v) = m^2c^4(u,v)$ since $\beta^2 = I$.
\end{proof}

Substituting \cref{eq:diracbilinear} into \cref{eq:qz},
\begin{equation}
  \label{eq:diracqz}
  q_z(u,v) = c^2(\nabla u, \nabla v) + m^2c^4(u,v)
    - \bar z\,(\Dop u, v) - z\,(u, \Dop v) + \abs{z}^2 (u,v),
\end{equation}
which is a bounded sesquilinear form on $H^1 \times H^1$ in which every term
carries at most one derivative.  The Riesz map has disappeared entirely.  So
\emph{standard $H^1$-conforming Lagrange elements are conforming for the folded
Dirac problem}, even though the folded operator $\Dop^*\Dop$ is second order.
This is \cref{sec:rayleigh} in miniature.  What controls conformity is the graph
space $\dom(\Dop) = H^1$, not the order of $\Dop^*\Dop$.

On a bounded domain the picture changes in the way \cref{sec:twodomains}
predicts.  Both integrations by parts in \cref{lem:diracform} acquire boundary
terms, $\dom(\Dop)$ is no longer all of $H^1$ but $H^1$ intersected with a
boundary condition of MIT bag type, and that condition has to be built into the
trial space as an essential condition.  \Cref{eq:diracqz} then carries an
additional boundary form.  Route (a) of \cref{sec:tworoutes} remains valid
verbatim and is what we would recommend using.

\subsection{Advection-diffusion, where it bites}
\label{sec:advdiff}

Take $\Lop u = -\nu\Delta u + \beta\cdot\nabla u$ on a bounded Lipschitz domain
$\Omega \subset \R^d$ with homogeneous Dirichlet conditions, $\Hb = L^2(\Omega)$
and $\Vb = H^1_0(\Omega)$, so that $a(u,v) = \nu(\nabla u, \nabla v) +
(\beta\cdot\nabla u, v)$.  This is the running example of
\cite[\S3]{LazzarinoNakatsukasaZerbinati2026} and it is genuinely non-normal,
which is the regime the whole exercise is aimed at.  Here
\begin{equation}
  \label{eq:advdiffdom}
  \dom(A) = H^2(\Omega) \cap H^1_0(\Omega) \subsetneq \Vb ,
\end{equation}
on a convex or $C^{1,1}$ domain, and the inclusion is strict.  Consequently
\begin{equation}
  \label{eq:advdiffqz}
  q_z(u,v) = \bigl(-\nu\Delta u + \beta\cdot\nabla u - zu,\
                    -\nu\Delta v + \beta\cdot\nabla v - zv\bigr)_{L^2}
\end{equation}
is a biharmonic-type form with form domain $H^2 \cap H^1_0$.  Continuous
piecewise linear elements are not in $H^2(\Omega)$, so they are not
$\dom(A)$-conforming, $\bigcup_h V_h$ is not a core of $A$, and neither the
one-sided bound \cref{eq:onesided} nor the convergence of \cref{prop:core} is
available.  Naively assembling \cref{eq:advdiffqz} on $\mathcal{P}_1$ Lagrange
elements is not a mildly suboptimal choice, it is outside the theory.

The honest options are the following, and we would like Matt's view on which he
intends.

\begin{itemize}
  \item Discretise $\dom(A)$ properly, with $C^1$ elements such as Argyris or
        Bogner--Fox--Schmit, or with $C^1$ splines.  This keeps the guarantee
        and is expensive.
  \item Rewrite the problem as a first-order system, for instance $\sigma =
        \nu\nabla u$ together with $-\divg\sigma + \beta\cdot\nabla u = f$, and
        fold that.  The system operator has graph space $H^1_0(\Omega) \times
        \Hdiv(\Omega)$ and standard elements are conforming.  But this computes
        the pseudospectra of a \emph{different} operator, acting on $L^2 \times
        (L^2)^d$.  Its spectrum agrees with that of $A$, its pseudospectra do
        not.  This is the same manoeuvre that \cite[\S3.4.4]{Colbrook2026}
        performs for Maxwell, where it is unobjectionable because the
        first-order system is the object of physical interest.  For
        advection-diffusion on $L^2(\Omega)$ it is a change of the question.
  \item Accept a non-conforming space and lose the certificate, keeping only
        the asymptotics.
\end{itemize}

It is worth noting that \cite{LazzarinoNakatsukasaZerbinati2026} chooses a
different Riesz map for exactly this reason.  Its $\tau$ is the one induced by
the $H^1_0$ inner product \cite[Equation~(3.20)]{LazzarinoNakatsukasaZerbinati2026},
and the resulting normal PDE
\begin{equation}
  \label{eq:normalpde}
  \nu(\nabla u, \nabla v)_{L^2}
    + \nu^{-1}\bigl(\Pi_{\nabla}(\beta u), \Pi_{\nabla}(\beta v)\bigr)_{L^2},
\end{equation}
from \cite[Equation~(3.27)]{LazzarinoNakatsukasaZerbinati2026}, is a perfectly
good $H^1_0$-conforming object.  Its Remark~3.7 warns that the $L^2$ Riesz map
``will not be consistent with the function spaces of the continuous problem''.
Read through \cref{sec:bridge}, that remark is exactly the statement
\cref{eq:advdiffdom}, that the $L^2$ Riesz map demands the graph space and not
the energy space.  \Cref{sec:pivot} takes up what this costs.

% ═════════════════════════════════════════════════════════════════════════════
\section{An affine structure in \texorpdfstring{$z$}{z}}
\label{sec:affine}

The outer layer of adaptivity samples $z$ over a region of the complex plane,
and it is worth recording that the $z$-dependence of the folded problem is
completely explicit.  Expanding \cref{eq:qz} on $\dom(A)$,
\begin{equation}
  \label{eq:zaffineform}
  q_z(u,v) = n(u,v) - \bar z\, a(u,v) - z\, a^{\adj}(u,v)
             + \abs{z}^2 m(u,v),
\end{equation}
where $n(u,v) := (Au,Av)_{\Hb}$, $a(u,v) = (Au,v)_{\Hb}$, $a^{\adj}(u,v) :=
\overline{a(v,u)} = (u,Av)_{\Hb}$ and $m(u,v) := (u,v)_{\Hb}$.  Equivalently, at
the level of operators,
\begin{equation}
  \label{eq:zaffineop}
  \Bop_z = \Aop^{\adj}\Riesz\Aop - z\,\Aop^{\adj} - \bar z\,\Aop
           + \abs{z}^2 \Mass ,
\end{equation}
using $\Riesz\Mass = \iota$ and $\Mass\Riesz\Mass = \Mass$ from \cref{eq:RM}.
In matrices this reads
\begin{equation}
  \label{eq:zaffinemat}
  \mathsf{L}^{(1)}(z) = \mathsf{N} - \bar z\,\mathsf{A} - z\,\mathsf{A}^{*}
    + \abs{z}^2\mathsf{M},
  \qquad
  \mathsf{N} = \mathsf{A}^{*}\mathsf{M}^{-1}\mathsf{A} ,
\end{equation}
which agrees with the expansion we had obtained by hand.  So four
$z$-independent matrices are assembled once and the whole complex-plane sweep is
a scalar recombination.  The factor $\mathsf{N}$ should be stored in the sparse
factored form \cref{eq:sparsefactor} rather than as a dense product.  This is
the same structural saving that \cite[Algorithm~3.5]{Colbrook2026} achieves at
the level of the square root, by precomputing a QR factorisation of the Gram
factor and reusing it across all sample points, and it appears to carry over to
the finite element setting unchanged.

% ═════════════════════════════════════════════════════════════════════════════
\section{How far can the pivot space be moved?}
\label{sec:pivot}

We can now say why the framework is built on an endomorphism of $\Hb$, rather
than on a map $\Vb \to \Vb'$, and how much freedom there is.

Suppose one tried to work with the triple directly and to measure everything in
$\Vb$, replacing \cref{eq:gamma} by
\begin{equation}
  \label{eq:gammaV}
  \gamma_{\Vb}(z) := \inf_{u \in \Vb\setminus\{0\}}
    \frac{\norm{(\Aop - z\Mass)u}_{\Vb'}}{\norm{u}_{\Vb}} .
\end{equation}
This is attractive because it is defined on all of $\Vb$, so the conformity
problem of \cref{sec:advdiff} evaporates and $\mathcal{P}_1$ elements are
admissible.  It is also, up to the choice of $\tau$, the quantity whose
conditioning governs the preconditioned normal equation of
\cite{LazzarinoNakatsukasaZerbinati2026}.
The difficulty is that the one-sidedness \cref{eq:onesided} is lost.  The dual
norm $\norm{g}_{\Vb'} = \sup_{v\in \Vb} \abs{\dual{g}{v}}/\norm{v}_{\Vb}$ is a
supremum, and restricting it to $V_h$ gives a \emph{lower} bound, while
restricting the outer infimum to $V_h$ gives an \emph{upper} bound.  The two
restrictions pull in opposite directions and the sign of the resulting error is
not determined.  Recovering a certificate then requires controlling the discrete
dual norm, which is the business of discontinuous Petrov--Galerkin methods with
optimal test functions \cite{DemkowiczGopalakrishnan2011}.

By contrast, when the target space is the pivot space, $\norm{g}_{\Hb'} =
\norm{\Riesz g}_{\Hb}$ is computed by an \emph{integral} rather than by a
supremum, and it is evaluated exactly by quadrature.  The only restriction is
then the outer one, and \cref{eq:onesided} survives.  That is the reason the
framework wants an endomorphism.  Its guarantee is a consequence of the residual
norm being exactly computable.

What survives of the freedom to choose an inner product is the following.

\begin{proposition}[reweighting the pivot]
  \label{prop:reweight}
  Let $W$ be bounded, self-adjoint and positive on $\Hb$ with bounded inverse,
  and let $(u,v)_W := (Wu,v)_{\Hb}$, so that $c\norm{u} \le \norm{u}_W \le
  C\norm{u}$ with $c^2 = \min\Sp(W)$ and $C^2 = \max\Sp(W)$.  Write $\kappa :=
  C/c$ and let $\gamma_W$ denote \cref{eq:gamma} computed in $\Hb_W := (\Hb,
  (\cdot,\cdot)_W)$.  Then, for every $z \in \C$,
  \begin{equation}
    \label{eq:reweight}
    \kappa^{-1}\gamma(z,A) \ \le\ \gamma_W(z,A) \ \le\ \kappa\,\gamma(z,A) .
  \end{equation}
  Consequently $\Sp(A)$ is unchanged, and for every $\epsilon > 0$
  \begin{equation}
    \label{eq:sandwich}
    \Sp_{\epsilon/\kappa}(A) \ \subseteq\ \Sp^W_\epsilon(A)
      \ \subseteq\ \Sp_{\kappa\epsilon}(A) .
  \end{equation}
\end{proposition}

\begin{proof}
  $\Hb_W$ is a Hilbert space with the same underlying vector space, so
  \cite[Lemma~1.2.5]{Colbrook2026} gives $\gamma_W(z,A) = \norm{(A -
  zI)^{-1}}_{\mathcal{L}(\Hb_W)}^{-1}$.  For any bounded $T$ on $\Hb$,
  $\norm{T}_{\mathcal{L}(\Hb_W)} = \norm{W^{1/2}TW^{-1/2}}_{\mathcal{L}(\Hb)}$,
  and $\norm{W^{1/2}}\,\norm{W^{-1/2}} = C/c = \kappa$, so
  $\kappa^{-1}\norm{T} \le \norm{T}_{\mathcal{L}(\Hb_W)} \le \kappa\norm{T}$.
  Apply this with $T = (A - zI)^{-1}$ and invert.  For \cref{eq:sandwich}, if
  $\gamma_W(z,A) < \epsilon$ then $\gamma(z,A) \le \kappa\gamma_W(z,A) <
  \kappa\epsilon$, and if $\gamma(z,A) < \epsilon/\kappa$ then $\gamma_W(z,A)
  \le \kappa\gamma(z,A) < \epsilon$.  Finally $\gamma_W = 0$ if and only if
  $\gamma = 0$.
\end{proof}

So the pivot inner product may be reweighted, the spectrum is untouched, and the
pseudospectra are sandwiched with explicit and computable constants.  If in
addition $W$ is local enough that $\norm{\cdot}_W$ is still evaluated exactly by
quadrature, which holds for multiplication operators and, discretely, for block
diagonal mass matrices on discontinuous spaces, then
\cite[Algorithm~3.4]{Colbrook2026} run in $\Hb_W$ still returns an upper bound
for $\gamma_W$, and \cref{eq:sandwich} converts that into a rigorous enclosure
for the original pseudospectra at a cost of one factor of $\kappa$.

This is the degree of freedom that
\cite{LazzarinoNakatsukasaZerbinati2026} exploits for preconditioning, and
\cref{prop:reweight} is what delimits the subset of those choices that preserve
Colbrook's certificate.  Choosing $W$ to improve the conditioning of the folded
pencil \cref{eq:pencil}, which squares the conditioning of $A - zI$, while
keeping $\kappa$ modest, looks to us like the most promising joint direction.
The relevant literature on choosing such inner products for operator
preconditioning is \cite{Kirby2010,MardalWinther2011}.

% ═════════════════════════════════════════════════════════════════════════════
\section{Questions}
\label{sec:questions}

\begin{question}[reweighting the pivot]
  Is \cref{prop:reweight} a route you have considered?  Trading a factor
  $\kappa$ in the reported $\epsilon$ against the conditioning of the folded
  pencil seems attractive, since folding squares the conditioning of $A - zI$.
  Does the SCI classification of \cite[Theorem~3.4.6]{Colbrook2026} survive the
  reweighting, given that the evaluation set would have to include the weights?
\end{question}

\begin{question}[mixed reformulation, same operator or a different one?]
  In \cref{sec:advdiff} the first-order system reformulation changes the
  operator whose pseudospectra are computed, even though the spectra agree.  In
  \cite[\S3.4.4]{Colbrook2026} the same manoeuvre is unobjectionable because the
  first-order Maxwell system is the object of interest.  Do you read the phrase
  ``reformulated into a mixed form using lower-order derivatives''
  \cite[p.~129]{Colbrook2026} as computing the same pseudospectra, or as a
  deliberate change of the operator?  For a second-order operator on $L^2$ we do
  not see how to avoid $C^1$ elements without giving up \cref{eq:onesided}.
\end{question}

\begin{question}[error control for non-normal problems]
  \Cref{rem:whereboffi} notes that \cite[Proposition~3.3.12 and
  Theorem~3.3.11]{Colbrook2026} both assume $A$ is normal, whereas the whole
  point of pseudospectra is the non-normal case.  The folded operator
  \emph{is} self-adjoint, so Kato--Temple applies to $\Bop_z$ provided the
  bottom of $\Sp(\Bop_z)$ is isolated and a lower bound on the gap is
  computable.  Is that the intended route to a posteriori control in the
  non-normal case?
\end{question}

\begin{question}[driving the inner adaptivity]
  When the bottom of $\Sp(\Bop_z)$ is essential spectrum, the standard finite
  element a posteriori theory for eigenvalues does not apply, and there is no
  gap to exploit.  Rayleigh--Ritz still certifies an upper bound, but it
  provides no local error indicator.  What would you use to drive mesh
  refinement in that regime?
\end{question}

% ═════════════════════════════════════════════════════════════════════════════
\bibliographystyle{plain}
\bibliography{references}

\end{document}
